\clearpage\pagestyle{otMain}\setlength{\otBlocEspacement}{2.5mm}\fontsize{11}{15}\selectfont\corrigesChapitre{1}{Les nombres jusqu’au million}
\begin{blocLivre}{cor-1-1}\begin{corrige}{1}{Écrire en chiffres}\demarche{Repérer la classe des mille dans le nom du nombre.}
On place chaque classe : $\res{8\,040}$, $\res{23\,405}$ et $\res{600\,009}$. Les rangs sans unités reçoivent un zéro.
\end{corrige}\end{blocLivre}
\begin{blocLivre}{cor-1-2}\begin{corrige}{2}{Lire en classes}\demarche{Séparer les classes avant de lire.}
On lit la classe des mille puis celle des unités : $\res{\text{quatre-vingt-six-mille-trois-cent-douze}}$ et $\res{\text{quatre-cent-cinquante-mille-soixante-dix}}$.
\end{corrige}\end{blocLivre}
\begin{blocLivre}{cor-1-3}\begin{corrige}{3}{La valeur d’un chiffre}\demarche{Associer chaque chiffre à son rang.}
Le chiffre des unités de mille est $\res2$. Le chiffre 8 occupe le rang des dizaines de mille : il vaut $\res{80\,000}$.
\end{corrige}\end{blocLivre}
\begin{blocLivre}{cor-1-4}\begin{corrige}{4}{Chiffre et nombre de dizaines}\demarche{Compter toutes les dizaines contenues.}
Le chiffre des dizaines est $\res2$. Comme $4\,627=462\times10+7$, le nombre de dizaines est $\res{462}$.
\end{corrige}\end{blocLivre}
\begin{blocLivre}{cor-1-5}\begin{corrige}{5}{Comparer deux nombres}\demarche{Comparer d’abord les longueurs, puis les rangs.}
Les centaines sont égales, mais $0<6$ aux dizaines : $\res{45\,806<45\,860}$. Un nombre à quatre chiffres est plus petit qu’un nombre à cinq chiffres : $\res{9\,875<12\,034}$.
\end{corrige}\end{blocLivre}
\begin{blocLivre}{cor-1-6}\begin{corrige}{6}{Décomposer}\demarche{N’écrire que les rangs non nuls.}
$304\,065=\res{3\times100\,000+4\times1\,000+6\times10+5}$. Les rangs vides n’apparaissent pas dans la somme.
\end{corrige}\end{blocLivre}
\begin{blocLivre}{cor-1-7}\begin{corrige}{7}{Recomposer}\demarche{Placer chaque valeur à son rang, compléter par des zéros.}
On place 7 aux dizaines de mille, 3 aux centaines, 9 aux unités : $\res{70\,309}$. Le second nombre est $\res{500\,020}$.
\end{corrige}\end{blocLivre}
\begin{blocLivre}{cor-1-8}\begin{corrige}{8}{Sur la demi-droite}\demarche{Chercher d’abord la valeur d’un intervalle.}
Deux graduations marquées voisines diffèrent de $1\,000$ ; chaque petit intervalle vaut donc 500. On lit $\res{A(25\,000)}$, $\res{B(26\,500)}$ et $\res{C(27\,500)}$.
\label{fig-cor-1-8}\figureGeo[corrige]{num-8}
\end{corrige}\end{blocLivre}
\begin{blocLivre}{cor-1-9}\begin{corrige}{9}{Encadrer}\demarche{Repérer le rang demandé avant d’encadrer.}
$\res{40\,000<45\,670<50\,000}$ : les dizaines de mille voisines sont $40\,000$ et $50\,000$. Pour les centaines : $\res{8\,200<8\,294<8\,300}$.
\end{corrige}\end{blocLivre}
\begin{blocLivre}{cor-1-10}\begin{corrige}{10}{Ordonner}\demarche{Grouper par nombre de chiffres puis comparer.}
Le nombre à quatre chiffres vient d’abord, celui à six chiffres en dernier. L’ordre est $\res{9\,980<55\,999<56\,040<56\,400<100\,000}$.
\end{corrige}\end{blocLivre}
\begin{blocLivre}{cor-1-11}\begin{corrige}{11}{Des fleuves à ranger}\demarche{Comparer les nombres avant de conclure sur les fleuves.}
L’ordre croissant des longueurs est $\res{647<775<812<1\,006}$ : Garonne, Seine, Rhône, Loire. La Garonne et la Seine mesurent moins de 800 km, car $\res{647<800}$ et $\res{775<800}$.
\end{corrige}\end{blocLivre}
\begin{blocLivre}{cor-1-12}\begin{corrige}{12}{Des carnets de timbres}\demarche{Relier carnets de dix et nombre de dizaines.}
$3\,240=324\times10$ : on remplit $\res{324}$ carnets, sans reste. $5\,278=527\times10+8$ : on remplit $\res{527}$ carnets et il reste $\res8$ timbres. Le nombre de carnets complets est le nombre de dizaines.
\end{corrige}\end{blocLivre}
\begin{blocLivre}{cor-1-13}\begin{corrige}{13}{Un nombre à inventer}\demarche{Fixer les rangs imposés, libérer les autres.}
Le nombre s’écrit $45\,37\square$ : par exemple $\res{45\,370}$. Le chiffre des unités est libre : il y a $\res{10}$ possibilités, de $45\,370$ à $45\,379$.
\end{corrige}\end{blocLivre}
\begin{blocLivre}{cor-1-14}\begin{corrige}{14}{Avec cinq cartes}\demarche{Choisir d’abord le chiffre de gauche.}
Pour le plus grand, ranger les cartes en ordre décroissant : $\res{97\,310}$. Pour le plus petit, commencer par le plus petit chiffre non nul puis ranger le reste en ordre croissant : $\res{10\,379}$. Un nombre à cinq chiffres ne commence pas par 0.
\end{corrige}\end{blocLivre}
\begin{blocLivre}{cor-1-15}\begin{corrige}{15}{Trois communes}\demarche{Traiter chaque question avec le même nombre bien identifié.}
L’ordre décroissant est $\res{45\,860>45\,806>9\,980}$. La plus grande population vérifie $\res{40\,000<45\,860<50\,000}$. Son chiffre des centaines est $\res8$.
\end{corrige}\end{blocLivre}
